\section{Architecture of the tokenizer}

The proposed direction is a learned tokenizer that compresses action chunks
while avoiding extreme concentrations of predictive difficulty. The original
draft does not yet specify its architecture; the description below outlines the
intended interface rather than a completed implementation.

During tokenizer training, an encoder maps an expert action chunk into a shorter
latent sequence. A discretization step converts that sequence into token IDs,
and a decoder reconstructs continuous actions. The exact encoder, quantizer,
decoder, and training losses remain to be supplied.

During policy training, those token IDs become prediction targets for one
autoregressive transformer conditioned on vision, language, and any robot-state
inputs. At inference, the transformer predicts action tokens and the action
decoder maps them into a continuous action chunk.

The action encoder processes demonstrations when creating training targets; it
does not require future expert actions at inference. Using a tokenizer and
action decoder alongside the policy is compatible with a single autoregressive
policy transformer, but the final architecture should make that boundary
explicit.

\begin{quote}
The design goal: preserve the motion that matters, with a token sequence the
policy can actually learn.
\end{quote}

Before presenting this as a result, specify the token budget, codebook or
vocabulary, action normalization, reconstruction objective, and mechanism
intended to balance information. Uniform token usage is not the same as evenly
distributed conditional predictive difficulty.
